% ============================================================
% 07-zakljucak.tex — PRVA VERZIJA
% ============================================================

U zaklju\v{c}ku dajemo osvrt na cjelokupni rad, sa\v{z}eto navodimo
\v{s}to je napravljeno i \v{s}to su klju\v{c}ni nalazi, a zatim
navodimo ograni\v{c}enja i otvaramo smjernice za budu\'{c}i rad.

Kroz rad je najprije obra\dj{}en teorijski okvir: problem
neuravnote\v{z}enosti klasa, temeljni SMOTE algoritam i pregled njegovih
jedanaest izvedenica. Zatim je opisano ostvareno programsko rje\v{s}enje
— Python implementacija svih varijanti, evaluacijski okvir s osam
klasifikatora i vi\v{s}e metrika, te modul za statisti\v{c}ku obradu
prema Dem\v{s}arovoj metodologiji. Provedena je eksperimentalna analiza
na stvarnim i sinteti\v{c}kim skupovima razli\v{c}itih karakteristika, a
rezultati su statisti\v{c}ki obra\dj{}eni i interpretirani. Kao dodatak,
izra\dj{}eno je web su\v{c}elje za vizualizaciju.

Klju\v{c}ni nalazi koje o\v{c}ekujemo: identificirat \'{c}emo koji se
SMOTE algoritmi najbolje pokazuju u pojedinim uvjetima, koliko
geometrijska pro\v{s}irenja zaista doprinose u odnosu na osnovni SMOTE,
te kakvu ulogu igra izbor klasifikatora.

Ograni\v{c}enja rada: fokus je na binarnoj klasifikaciji, \v{s}to
zna\v{c}i da zaklju\v{c}ci nisu izravno primjenjivi na vi\v{s}eklasne
probleme. Tako\dj{}er, broj testiranih vrijednosti parametara je
ograni\v{c}en, a sinteti\v{c}ki skupovi ne mogu u potpunosti zamijeniti
stvarne podatke.

Za budu\'{c}i rad otvaramo nekoliko pravaca: pro\v{s}irenje na
vi\v{s}eklasne probleme, dublja integracija s deep learning pristupima,
automatska optimizacija hiperparametara SMOTE algoritama te evaluacija
na ekstremno neuravnote\v{z}enim skupovima (IR $>$ 100) koji su u praksi
sve \v{c}e\v{s}\'{c}i.
